%%%%%%%%%%%%%%%%%%%%%%%%%%%%%%%%%%%%%%%%%
% Programming/Coding Assignment
% LaTeX Template
%
% This template has been downloaded from:
% http://www.latextemplates.com
%
% Original author:
% Ted Pavlic (http://www.tedpavlic.com)
%
% Note:
% The \lipsum[#] commands throughout this template generate dummy text
% to fill the template out. These commands should all be removed when 
% writing assignment content.
%
% This template uses a Perl script as an example snippet of code, most other
% languages are also usable. Configure them in the "CODE INCLUSION 
% CONFIGURATION" section.
%
%%%%%%%%%%%%%%%%%%%%%%%%%%%%%%%%%%%%%%%%%

%----------------------------------------------------------------------------------------
%	PACKAGES AND OTHER DOCUMENT CONFIGURATIONS
%----------------------------------------------------------------------------------------

\documentclass{article}
\usepackage[utf8]{inputenc}
\usepackage{listingsutf8}
\usepackage[spanish]{babel}




\usepackage{fancyhdr} % Required for custom headers
\usepackage{lastpage} % Required to determine the last page for the footer
\usepackage{extramarks} % Required for headers and footers
\usepackage[usenames,dvipsnames]{color} % Required for custom colors
\usepackage{graphicx} % Required to insert images
\usepackage{listings} % Required for insertion of code
\usepackage{courier} % Required for the courier font
\usepackage{lipsum} % Used for inserting dummy 'Lorem ipsum' text into the template
\usepackage{svg}
\usepackage{attachfile}
\usepackage{currfile}
\usepackage{multicol}
\usepackage{alltt}
\usepackage{framed}

\hypersetup{
colorlinks,
citecolor=black,
filecolor=black,
linkcolor=black,
urlcolor=blue
}


% Margins
\topmargin=-0.45in
\evensidemargin=0in
\oddsidemargin=0in
\textwidth=6.5in
\textheight=9.0in
\headsep=0.25in

\linespread{1.1} % Line spacing

% Set up the header and footer
\pagestyle{fancy}
%\lhead{\hmwkAuthorName} % Top left header
\lhead{\hmwkClass}
\rhead{\hmwkTitle} % Top right header
\chead{} % Top center head
\lfoot{\lastxmark} % Bottom left footer
\cfoot{} % Bottom center footer
\rfoot{ \thepage\ / \protect\pageref{LastPage}} % Bottom right footer
\renewcommand\headrulewidth{0.4pt} % Size of the header rule
\renewcommand\footrulewidth{0.4pt} % Size of the footer rule

\setlength\parindent{0pt} % Removes all indentation from paragraphs

%----------------------------------------------------------------------------------------
%	CODE INCLUSION CONFIGURATION
%----------------------------------------------------------------------------------------

\definecolor{MyDarkGreen}{rgb}{0.0,0.4,0.0} % This is the color used for comments
\lstloadlanguages{Perl} % Load Perl syntax for listings, for a list of other languages supported see: ftp://ftp.tex.ac.uk/tex-archive/macros/latex/contrib/listings/listings.pdf
\lstset{language=Perl, % Use Perl in this example
frame=single, % Single frame around code
basicstyle=\small\ttfamily, % Use small true type font
keywordstyle=[1]\color{Blue}\bf, % Perl functions bold and blue
keywordstyle=[2]\color{Purple}, % Perl function arguments purple
keywordstyle=[3]\color{Blue}\underbar, % Custom functions underlined and blue
identifierstyle=, % Nothing special about identifiers                                         
commentstyle=\usefont{T1}{pcr}{m}{sl}\color{MyDarkGreen}\small, % Comments small dark green courier font
stringstyle=\color{Purple}, % Strings are purple
showstringspaces=false, % Don't put marks in string spaces
tabsize=5, % 5 spaces per tab
%
% Put standard Perl functions not included in the default language here
morekeywords={rand},
%
% Put Perl function parameters here
morekeywords=[2]{on, off, interp},
%
% Put user defined functions here
morekeywords=[3]{test},
%
morecomment=[l][\color{Blue}]{...}, % Line continuation (...) like blue comment
numbers=left, % Line numbers on left
firstnumber=1, % Line numbers start with line 1
numberstyle=\tiny\color{Blue}, % Line numbers are blue and small
stepnumber=5 % Line numbers go in steps of 5
}

% Creates a new command to include a perl script, the first parameter is the filename of the script (without .pl), the second parameter is the caption
\newcommand{\perlscript}[2]{
\begin{itemize}
\item[]\lstinputlisting[caption=#2,label=#1]{#1.pl}
\end{itemize}
}

%----------------------------------------------------------------------------------------
%	DOCUMENT STRUCTURE COMMANDS
%	Skip this unless you know what you're doing
%----------------------------------------------------------------------------------------

% Header and footer for when a page split occurs within a problem environment
\newcommand{\enterProblemHeader}[1]{
%\nobreak\extramarks{#1}{#1 continued on next page\ldots}\nobreak
%\nobreak\extramarks{#1 (continued)}{#1 continued on next page\ldots}\nobreak
}

% Header and footer for when a page split occurs between problem environments
\newcommand{\exitProblemHeader}[1]{
%\nobreak\extramarks{#1 (continued)}{#1 continued on next page\ldots}\nobreak
%\nobreak\extramarks{#1}{}\nobreak
}

\setcounter{secnumdepth}{0} % Removes default section numbers
\newcounter{homeworkProblemCounter} % Creates a counter to keep track of the number of problems

\newcommand{\homeworkProblemName}{}
\newenvironment{homeworkProblem}[1][]{ % Makes a new environment called homeworkProblem which takes 1 argument (custom name) but the default is "Problem #"
\stepcounter{homeworkProblemCounter} % Increase counter for number of problems
\renewcommand{\homeworkProblemName}{Ejercicio \arabic{homeworkProblemCounter} #1} % Assign \homeworkProblemName the name of the problem
\section{\homeworkProblemName} % Make a section in the document with the custom problem count
\enterProblemHeader{\homeworkProblemName} % Header and footer within the environment
}{
\exitProblemHeader{\homeworkProblemName} % Header and footer after the environment
%\clearpage
}


\newcommand{\problemAnswer}[1]{ % Defines the problem answer command with the content as the only argument
\noindent\framebox[\columnwidth][c]{\begin{minipage}{0.98\columnwidth}#1\end{minipage}} % Makes the box around the problem answer and puts the content inside
}

\newcommand{\homeworkSectionName}{}
\newenvironment{homeworkSection}[1]{ % New environment for sections within homework problems, takes 1 argument - the name of the section
\renewcommand{\homeworkSectionName}{#1} % Assign \homeworkSectionName to the name of the section from the environment argument
\subsection{\homeworkSectionName} % Make a subsection with the custom name of the subsection
\enterProblemHeader{\homeworkProblemName\ [\homeworkSectionName]} % Header and footer within the environment
}{
\enterProblemHeader{\homeworkProblemName} % Header and footer after the environment
}

%----------------------------------------------------------------------------------------
%	NAME AND CLASS SECTION
%----------------------------------------------------------------------------------------

\newcommand{\hmwkTitle}{Redefine el hmwkTitle} % Assignment title
\newcommand{\hmwkDueDate}{asdfadsf} % Due date
\newcommand{\hmwkClass}{Seguridad informática} % Course/class
\newcommand{\hmwkClassTime}{} % Class/lecture time
\newcommand{\hmwkClassInstructor}{} % Teacher/lecturer
\newcommand{\hmwkAuthorName}{Álvaro González Sotillo} % Your name

%----------------------------------------------------------------------------------------
%	TITLE PAGE
%----------------------------------------------------------------------------------------

\title{
\vspace{2in}
\textmd{\textbf{\hmwkClass:\ \hmwkTitle}}\\
\vspace{0.1in}\large{\textit{\hmwkClassInstructor\ \hmwkClassTime}}
\vspace{3in}
}

\author{\textbf{\hmwkAuthorName}}
\date{} % Insert date here if you want it to appear below your name


%----------------------------------------------------------------------------------------

\usepackage{fancybox}
\newcommand{\codigo}[1]{\texttt{#1}}


% CUADRITO
\newsavebox{\fmboxx}
\newenvironment{cuadrito}[1][14cm]
{\noindent \begin{center} \begin{lrbox}{\fmboxx}\noindent\begin{minipage}{#1}}
{\end{minipage}\end{lrbox}\noindent\shadowbox{\usebox{\fmboxx}} \end{center}}





\newenvironment{entradasalida}[2][14cm]
{
  \newcommand{\elnombredelafiguradeentradasalida}{#2}
  \begin{figure}[h]
    \begin{cuadrito}[#1]
      \begin{scriptsize}
\begin{alltt}
}
{%
\end{alltt}%
      \end{scriptsize}%
    \end{cuadrito}%
    \caption{\elnombredelafiguradeentradasalida}
  \end{figure}
}


\newenvironment{entradasalidacols}[2][14cm]
{
\newcommand{\elnombredelafiguradeentradasalida}{#2}
%\begin{wrapfigure}{}{0.1}
\begin{cuadrito}[#1]
\begin{scriptsize}
\begin{alltt}
}
{
\end{alltt}
\end{scriptsize}
\end{cuadrito}
\captionof{figure}{\elnombredelafiguradeentradasalida}
%\end{wrapfigure}
}


\newcommand{\ficheroautoref}[0]{

\attachfile[mimetype=text/plain,
description={La plantilla},
subject={La plantilla}]
{../common/plantilla-ejercicio.tex}

\attachfile[mimetype=text/plain,
description={El fichero TEX original para crear este documento, no sea que se nos pierda},
subject={El fichero TEX original para crear este documento, no sea que se nos pierda}]
{\currfilename}}

\newcommand{\entradausuario}[1]{\textit{\textbf{#1}}}

\newcommand{\enlace}[2]{\textcolor{blue}{{\href{#1}{#2}}}}

\newcommand{\adjuntarfichero}[3]{
\textattachfile[mimetype=text/plain,
color={0 0 0},
description={#3},
subject={#1}]
{#1}
{\textcolor{blue}{\codigo{#2}}}
}

\newcommand{\adjuntardoc}[2]{
\textattachfile[mimetype=text/plain,
color={0 0 0},
description={#1},
subject={#1}]
{#1}
{\textcolor{blue}{#2}}
}


\newcommand{\plantilladeclase}[2]{
\adjuntarfichero{#1.java}{#2}{Plantilla para la clase #2}
}

\lstset{
  literate=%
    {á}{{\'a}}1
    {é}{{\'e}}1
    {í}{{\'i}}1
    {ó}{{\'o}}1
    {ú}{{\'u}}1
    {Á}{{\'A}}1
    {É}{{\'E}}1
    {Í}{{\'Y}}1
    {Ó}{{\'O}}1
    {Ú}{{\'U}}1
}

\renewcommand{\lstlistingname}{Listado}
\captionsetup[lstlisting]{font={footnotesize},labelfont=bf,position=bottom}
\captionsetup[figure]{font={footnotesize},labelfont=bf}
\lstnewenvironment{listadojava}[1]
{
  \lstset{captionpos=b,caption={#1},language=Java,frame=single,basicstyle=\footnotesize\ttfamily,showstringspaces=false,numbers=left,xleftmargin=2em}
}
{
}

% LISTADO SHELL
\lstnewenvironment{listadoshell}[2][]
{
  \lstset{captionpos=b,caption={#2},label={#1},language=bash,frame=single,basicstyle=\scriptsize\ttfamily,showstringspaces=false,numbers=none}
}
{
}

% LISTADO TXT
\lstnewenvironment{listadotxt}[2][]
{
  \lstset{inputencoding=utf8,captionpos=b,caption={#2},label={#1},frame=single,basicstyle=\scriptsize\ttfamily,showstringspaces=false,numbers=none}
}
{
}


\newcommand{\graficosvg}[3][14cm]{
\begin{figure}[htbp]
\centering
\textattachfile{#2.svg}{
\color{black}
\includesvg[width=#1]{#2}
}
\caption{#3}
\end{figure}
}


\newcommand{\graficosvguml}[3][7cm]{
  \texttt{\graficosvg[#1]{#2}{#3}}
}


\newcommand{\primerapagina}{
\newpage
\tableofcontents
{\ficheroautoref}
\newpage
}

\usepackage{eurosym}
\usepackage{qrcode}
\usepackage{ragged2e}



\renewcommand{\hmwkTitle}{Práctica de seguridad en el sistema}


\usepackage{blindtext}

\newcommand{\dibujoadmin}{\raisebox{-2mm}{\includegraphics[height=7mm]{admin.pdf}}}
\newcommand{\dibujohacker}{\raisebox{-2mm}{\includegraphics[height=7mm]{hacker.pdf}}}


% CAJAS
\usepackage[skins]{tcolorbox}
\newtcolorbox{aActividad}[1][Actividad]{
  enhanced,
  colback=gray!5!white,
  colframe=gray!75!black,fonttitle=\bfseries,
  colbacktitle=gray!85!black,
  attach boxed title to top left={yshift=-2mm,xshift=2mm},
  title=#1
}

\newtcolorbox{ActividadHacker}[1][Actividad]{
  enhanced,
  colback=gray!5!white,
  colframe=gray!75!black,fonttitle=\bfseries,
  colbacktitle=gray!85!black,
  attach boxed title to top left={yshift=-2mm,xshift=2mm},
  title=#1\hspace{2mm}\dibujohacker
}

\newtcolorbox{ActividadAdmin}[1][Actividad]{
  enhanced,
  colback=gray!5!white,
  colframe=gray!75!black,fonttitle=\bfseries,
  colbacktitle=gray!85!black,
  attach boxed title to top left={yshift=-2mm,xshift=2mm},
  title=#1\hspace{1mm}\dibujoadmin
}

\usepackage{multicol}

\begin{document}

% \maketitle

% ----------------------------------------------------------------------------------------
%	TABLE OF CONTENTS
% ----------------------------------------------------------------------------------------

% \setcounter{tocdepth}{1} % Uncomment this line if you don't want subsections listed in the ToC

{\primerapagina}

\setlength{\parindent}{0em}
\setlength{\parskip}{1em}


\section{Objetivo de la práctica}
Durante el desarrollo de esta práctica, el alumno provocará varios agujeros de seguridad en sistemas Linux, que un usuario sin privilegios podría aprovechar para tomar el control del sistema. De esta forma podrá comprobar como errores más o menos sutiles pueden convertir un sistema en inseguro. Además, investigará qué tipo de niveles de acceso puede tener un usario en el sistema y la forma de modificarlos.

Cada paso está planteado desde el punto de vista de:

\begin{tabular}{r c p{10cm}}
  %\hline
  Un administrador de sistema & \colorbox{gray!85}{\dibujoadmin} & Intentará detectar, evitar o paliar los ataques a la seguridad. También cometerá errores que dejarán el sistema vulnerable. \\
  % \hline
   Un \textit{hacker} & \colorbox{gray!85}{\dibujohacker} & Intentará aprovechar fallos en la seguridad \\
  %\hline
\end{tabular}


\begin{homeworkProblem}[: Aprovecha un fallo de seguridad en \textbf{Linux} (2 puntos)]

  \begin{multicols}{2}

    \null \vfill
    \begin{ActividadAdmin}[1. Problema de seguridad]
      ¡El usuario \texttt{root} a dejado a \texttt{777} los permisos de \texttt{/etc/groups} y \texttt{/etc/passwd}!


    \end{ActividadAdmin}
    \vfill \null
    \columnbreak



    \begin{ActividadHacker}[2. Un \textit{hacker} lo aprovecha]
      Tienes acceso al usuario \texttt{normal}, pero no tiene ningún permiso especial. ¡Aprovecha el problema de seguridad y consigue ser \texttt{root}!
      \begin{itemize}
      \item Puedes cambiar la password del \texttt{root}, o bien…
      \item Puedes añadir usuarios al grupo \texttt{admin} y al grupo \texttt{sudo}, que son especiales para el comando \texttt{sudo}
      \end{itemize}
    \end{ActividadHacker}



  \end{multicols}


\end{homeworkProblem}


\begin{homeworkProblem}[: Elección de una buena contraseña (1 punto)]

  \begin{multicols}{2}

    \null \vfill
    \begin{ActividadAdmin}[1. Problema de seguridad]
      Un usuario utiliza una contraseña muy débil (\textbf{pass}). Su \textit{hash} se guarda en el fichero \path{/etc/shadow},
    \end{ActividadAdmin}
    \vfill \null
    \columnbreak


    \begin{ActividadHacker}[2. Un \textit{hacker} lo aprovecha]
      \begin{itemize}
      \item Consigues de un viejo \textit{backup} el fichero \path{/etc/shadow}, así que conoces el \textit{hash} md5 de la contraseña (en este caso, calcula el \textit{hash} con el comando \texttt{md5sum}).  
      \item Intenta encontrar el \textit{hash} de la contraseña en Internet. Si lo encuentras, es que la contraseña \textbf{pass} no es demasiado buena, y cualquiera puede encontrar la contraseña a partir del \textit{hash} 
      \end{itemize}
    \end{ActividadHacker}
  \end{multicols}

\end{homeworkProblem}

\newpage
\begin{homeworkProblem}[: Auditoría de accesos en \textbf{Linux} (1 punto)]

  \begin{multicols}{2}

    \begin{ActividadHacker}[1. Problema de seguridad]
      Alguien está intentando acceder a tu \textbf{Linux}, probando usuarios y contraseñas por fuerza bruta
      \begin{itemize}
      \item Intenta hacer \textit{login} con un usuario que no existe (por \textit{ssh} o en una consola de texto)
      \item Intenta hacer \texttt{sudo} con contraseña incorrecta
      \item Intenta hacer \texttt{sudo} con un usuario que no tiene acceso a \texttt{sudo}
      \end{itemize}
    \end{ActividadHacker}
    \columnbreak

    \null \vfill

    \begin{ActividadAdmin}[2. El administrador lo detecta]
      Como tienes sospechas de lo que está ocurriendo, decides consultar los ficheros de log (en \texttt{/var/log}). ¿En qué ficheros se guardan estos sucesos?. ¿Qué información se obtiene?
    \end{ActividadAdmin}

    \vfill \null

  \end{multicols}


\end{homeworkProblem}



\begin{homeworkProblem}[: Escalada de privilegios en \textbf{Linux} (2 puntos)]


  \begin{multicols}{2}

    \null \vfill
    \begin{ActividadAdmin}[1. Problema de seguridad]
      ¡El usuario \texttt{root} se ha dejado una consola abierta!.

      Ahora que lo has visto, quieres aprovechar la posibilidad para poder ser \texttt{root} otro día. Tu problema es que si cambias la contraseña \textbf{el administrador lo notaría}. Necesitas una forma que no afecte a la configuración del sistema (los ficheros de \path{/etc}).
    \end{ActividadAdmin}
    \vfill \null
    \columnbreak



    \begin{ActividadHacker}[2. El \textit{hacker} lo aprovecha]
      \begin{itemize}
      \item Consulta información del permiso \texttt{s} en el manual de \texttt{chmod}
      \item Crea el usuario \texttt{normal2}
      \item Crea una copia de la \textit{shell} (\texttt{/usr/bash}) en el directorio \textit{home} de \textbf{normal2} (\path{/home/normal2/shell-de-root})
      \item Haz que \texttt{shell-de-root} tenga como propietario a \texttt{root} y tenga el permiso \texttt{s}
      \item Ejecuta \texttt{shell-de-root -p} siendo el usuario \texttt{normal2}. ¿Qué ocurre?
      \end{itemize}
    \end{ActividadHacker}

  \end{multicols}

  
\end{homeworkProblem}


\newpage
\begin{homeworkProblem}[: Implementa el uso de cuotas de disco en \textbf{Linux} (3 puntos)]

  \begin{multicols}{2}

    \null \vfill
    \begin{ActividadHacker}[1. Problema de seguridad]
      El usuario \texttt{normal} está abusando del sistema, ya que se está bajando multitud de archivos de Internet, y no deja sitio libre para el trabajo legítimo de los demás usuarios (es un ataque \textit{Denial of Serivice}).
    \end{ActividadHacker}
    \vfill \null
    \columnbreak


    \begin{ActividadAdmin}[2. El administrador lo soluciona]
      \newcommand{\tutorialcuota}{https://www.linuxtotal.com.mx/?cont=info\_admon\_018}%{https://www.howtoforge.com/tutorial/linux-quota-ubuntu-debian}
      \begin{itemize}
      \item Consulta \enlace{\tutorialcuota}{este enlace} para ver cómo manejar las cuotas

        \centering\qrcode[height=4cm]{\tutorialcuota}
        
      \item \justify Haz que el usuario \texttt{normal2} tenga una cuota de 1000 KBytes, Observa qué ocurre cuando intenta utilizar más de este espacio del disco.
      \end{itemize}


    \end{ActividadAdmin}

  \end{multicols}



\end{homeworkProblem}


\begin{homeworkProblem}[: Listas de control de acceso (1 punto)]

  \begin{multicols}{2}

    \null \vfill
    \begin{ActividadAdmin}[1. Problema de seguridad]
      El administrador tiene ahora un ordenador con \textbf{Windows}. Para configurarlo adecuadamente, necesita  una lista de qué permisos pueden otorgarse/negarse en \textbf{Windows} a cada tipo de \textit{objeto}.
    \end{ActividadAdmin}
    \vfill \null
    \columnbreak



    \begin{ActividadAdmin}[2. El administrador lo soluciona]
      Consigue una lista de los permisos asignables a:
      \begin{itemize}
      \item Ficheros
      \item Procesos (utiliza el programa \texttt{procexp})
      \item Entradas de registro (utiliza el programa \texttt{regedit})
      \end{itemize}
    \end{ActividadAdmin}

  \end{multicols}

  

\end{homeworkProblem}



\section{Qué se valorará}
El resultado de la práctica debe ser una memoria con los pasos que el alumno ha seguido para seguir los pasos del administrador y el \textit{hacker}. Se valorará:
\begin{itemize}
\item Que cada paso quede bien documentado. 
\item La corrección técnica (que funcione)
\item Que esté correctamente redactado, de forma que nuestro lector lo entienda 
\item La apariencia profesional:
  \begin{itemize}
  \item Estética
  \item Organización
  \item Homogeneidad de formatos y estilos
  \end{itemize}
\end{itemize}

\section{Instrucciones de entrega}
\begin{itemize}
\item El ejercicio se realizará y entregará de manera individual.
  \begin{itemize}
  \item Solo se admiten trabajos en pareja, si en clase es necesario compartir ordenador.
  \item En este caso, todos los integrantes del grupo deben subir el trabajo al aula virtual.

  \end{itemize}
\item Los trabajos pueden entregarse:
  \begin{itemize}
  \item En formato \textbf{DOC} o \textbf{DOCX}.
  \item En formato \textbf{ODT}.
  \item En formato \textbf{PDF}. 
  \item Como una entrada en un \textbf{blog}

  \end{itemize}
  
\item La entrega se realizará en la tarea correspondiente del aula virtual. Si se entrega un fichero, este se subirá directamente. Si es una entrada de blog, se subirá un fichero de texto con la URL de dicha entrada.
\end{itemize}  

\section{Créditos}
\begin{itemize}
\item Icons made by \url{http://www.freepik.com} from \url{http://www.flaticon.com} are licensed by \textbf{CC 3.0 BY}
\item Icons made by Picol from \url{http://www.flaticon.com} are licensed by \textbf{CC 3.0 BY}
\end{itemize}
\end{document}
